We describe the three components of SpecSSM: (1)~the guided SSM drafter architecture, (2)~the generation loop with the activation replay resynchronization method, and (3)~CPU offloading of the drafter.
Figure~\ref{fig:architecture} provides an overview of the complete system.

\begin{figure}[t]
  \centering
  % TODO: Replace with hand-drawn/TikZ architecture diagram
  \fbox{\parbox{0.9\textwidth}{\centering\vspace{2em}
  \textbf{Architecture Overview (to be drawn)}\\[0.5em]
  \small
  Left: Verifier (LLaMA on GPU) $\to$ GuidanceExtractor (concat layers 5, 16, 29) $\to$ PrepMambaDeltas $\to$ per-layer $\Delta_\ell$\\
  Center: Mamba2 drafter (on CPU) with x-branch delta injection at each layer\\
  Right: Speculative decoding loop: Draft $K$ tokens $\to$ Verify in batch $\to$ Accept/Reject $\to$ Activation Replay\\
  Bottom: CPU-GPU pipeline timeline (9\,ms draft $\subset$ 16\,ms verify window)
  \vspace{2em}}}
  \caption{SpecSSM architecture. The verifier extracts hidden states from selected layers, which are projected into per-layer additive deltas and injected into the Mamba2 drafter's x-branch. After speculative rejection, activation replay resynchronizes the SSM state in $O(K)$ time. The drafter runs on CPU with a custom INT8 kernel while the verifier occupies the GPU.}
  \label{fig:architecture}
\end{figure}

\subsection{Guided SSM Drafter}
\label{sec:guided}

Our drafter is a Mamba2-27M model~\citep{dao2024transformers} with 16 layers, hidden size 512, expand factor 2 ($d_{\mathrm{inner}}=1024$), state size 128, 16 heads, and head dimension 64.
The 27M refers to the backbone (SSM layers + projections); the model additionally shares the verifier's 66M embedding/LM-head layer, making total parameter count tokenizer-dependent.
We modify the Mamba2 architecture to accept \emph{guidance signals} from the verifier.

\paragraph{Guidance extraction.}
Given a verifier $p_\theta$ (LLaMA 3.1-8B), we extract hidden states from three layers at positions early ($\ell_1=5$), middle ($\ell_2=16$), and late ($\ell_3=29$) in the verifier's 32-layer stack.
These are concatenated and projected into a single guidance embedding:
\begin{equation}
  \mathbf{g}_t = \mathbf{W}_g\, [\mathbf{h}_t^{(\ell_1)} \| \mathbf{h}_t^{(\ell_2)} \| \mathbf{h}_t^{(\ell_3)}] + \mathbf{b}_g,
  \label{eq:guidance_extract}
\end{equation}
where $\mathbf{W}_g \in \mathbb{R}^{d_v \times 3d_v}$, $d_v = 4096$ is the verifier hidden dimension, and $\mathbf{g}_t \in \mathbb{R}^{d_v}$.
The projection is initialized as the identity average: $\mathbf{W}_g = [\mathbf{I}\, \mathbf{I}\, \mathbf{I}]$, $\mathbf{b}_g = \mathbf{0}$.

\paragraph{Guidance injection via x-branch deltas.}
The guidance embedding is mapped to per-layer additive deltas through a learned linear projection:
\begin{equation}
  \Delta_\ell = (\mathbf{W}_{\mathrm{prep}} \cdot \mathrm{LN}(\mathbf{g}_t))_\ell, \quad \ell = 1, \ldots, L_d,
  \label{eq:prep_deltas}
\end{equation}
where $L_d$ is the number of drafter layers and $\Delta_\ell \in \mathbb{R}^{d_{\mathrm{inner}}}$.
At each drafter layer, after the \texttt{in\_proj} linear maps the input into the gate (z-branch) and state-update (x-branch) components, we inject the delta additively into the x-branch:
\begin{equation}
  \mathbf{x}_\ell \leftarrow \mathbf{x}_\ell + \Delta_\ell.
  \label{eq:inject}
\end{equation}
This targets the state-update pathway (the input to the selective scan), allowing guidance to influence the SSM's recurrent dynamics without altering the gating mechanism.
The projection weights are zero-initialized, so at the start of training the guided model behaves identically to the unguided base.

\paragraph{Training.}
The drafter is trained via knowledge distillation from the verifier using total variation distance (TVD) loss:
\begin{equation}
  \mathcal{L}_{\mathrm{TVD}} = \frac{1}{2} \sum_v |p_\theta(v \mid \mathbf{x}_{<t}) - q_\phi(v \mid \mathbf{x}_{<t}, \mathbf{g}_{<t})|,
  \label{eq:tvd}
\end{equation}
where $q_\phi$ denotes the guided drafter.
We train for 10 epochs on 50K filtered samples from UltraChat~\citep{ding2023ultrachat} using AdamW with learning rate $3{\times}10^{-4}$ and cosine decay.
The drafter backbone is optionally finetuned alongside the guidance projections; we find that backbone finetuning provides $3{\times}$ more benefit than frozen guidance alone (\S\ref{sec:analysis}).

\subsection{Speculative Decoding with SSM Cache Resynchronization}
\label{sec:replay}

\paragraph{Non-compact layout.}
We use a non-compact layout where the write cursor advances by $K+1$ per round regardless of how many tokens are accepted.
Rejected draft tokens remain in the buffer at their physical positions; a 2D attention mask zeros out rejected positions so the verifier never attends to stale key-value entries.
Position IDs encode \emph{semantic} positions, skipping over rejected slots.

\paragraph{The SSM cache problem.}
In transformer-based speculative decoding, rejection is handled by truncating the KV cache to the last accepted position.
For SSMs, the hidden state $\mathbf{h}_t$ encodes the entire history through the recurrent dynamics of Eq.~\eqref{eq:ssm}, and there is no efficient truncation operation.
After rejection at position $t$, the hidden state includes information from rejected tokens $t+1, \ldots, t+K$ in the convolution buffer and SSM state.

\paragraph{Activation replay.}
Our key insight is that instead of re-running the drafter over the entire prompt to reconstruct a clean state, we can \emph{replay} only the accepted tokens through the SSM.
Specifically, after the verifier accepts $n_a$ of $K$ draft tokens, we:
\begin{enumerate}
  \item Extract the accepted token subsequence $\hat{\mathbf{x}} = [x_{t+1}, \ldots, x_{t+n_a}, x_{\mathrm{resample}}]$
  \item Restore the SSM state $\mathbf{h}_t$ saved \emph{before} the current draft round began (snapshot cost: 8.96\,MB for Mamba2-27M at batch size 1)
  \item Replay $\hat{\mathbf{x}}$ through the drafter starting from $\mathbf{h}_t$
\end{enumerate}
This produces a state that is \emph{bit-identical} to what would result from never having drafted the rejected tokens, in $O(n_a)$ time rather than $O(T)$ for full re-prefill (where $T$ is the total sequence length).
On CPU, where single-step latency is $\sim$7\,ms, this difference is dramatic: at 1024-token context, re-prefill costs 8.0\,s while replay costs 40\,ms ($201\times$ speedup, Table~\ref{tab:replay}).
On GPU the benefit is less pronounced since the model is small enough for fast parallel re-prefill ($\sim$16\,ms regardless of context length).

Algorithm~\ref{alg:specssm} summarizes the full SpecSSM generation loop with activation replay.

\begin{algorithm}[t]
\caption{SpecSSM: Speculative decoding with guided SSM drafter and activation replay.}
\label{alg:specssm}
\begin{algorithmic}[1]
\REQUIRE Verifier $p_\theta$ (GPU), guided drafter $q_\phi$ (CPU), prompt $\mathbf{x}_{1:n}$, lookahead $K$
\STATE Prefill: run $p_\theta(\mathbf{x}_{1:n})$, extract guidance $\mathbf{g} \gets \textsc{ExtractGuidance}(p_\theta)$
\STATE Prefill: run $q_\phi(\mathbf{x}_{1:n-1})$ with zero deltas to warm SSM state $\mathbf{h}$
\WHILE{not done}
  \STATE $\mathbf{h}_{\mathrm{snap}} \gets \mathbf{h}$ \COMMENT{Snapshot SSM state (8.96\,MB)}
  \STATE $\boldsymbol{\Delta} \gets \textsc{PrepDeltas}(\mathbf{g})$ \COMMENT{Project guidance to per-layer deltas}
  \FOR{$k = 1$ \TO $K$}
    \STATE $\hat{x}_k \sim q_\phi(\cdot \mid \mathbf{h}, \boldsymbol{\Delta})$ \COMMENT{Draft one token with guidance (CPU)}
  \ENDFOR
  \STATE $\mathbf{p}_{1:K+1} \gets p_\theta(\hat{x}_{1:K} \mid \text{mask})$ \COMMENT{Verify all drafts in one pass (GPU)}
  \STATE $n_a \gets \textsc{RejectSample}(\mathbf{p}, \mathbf{q}, \hat{x}_{1:K})$ \COMMENT{Accept prefix, resample at rejection}
  \STATE $\mathbf{h} \gets \mathbf{h}_{\mathrm{snap}}$ \COMMENT{Restore clean state}
  \STATE Replay: run $q_\phi([x_{\mathrm{accepted}_1}, \ldots, x_{\mathrm{accepted}_{n_a}}, x_{\mathrm{resample}}])$ from $\mathbf{h}$ \COMMENT{$O(n_a)$ replay}
  \STATE $\mathbf{g} \gets \textsc{ExtractGuidance}(p_\theta, \text{position}=n_a)$ \COMMENT{Fresh guidance}
  \STATE Update cursor, zero rejected positions in attention mask
\ENDWHILE
\end{algorithmic}
\end{algorithm}

\subsection{CPU Offloading}
\label{sec:cpu}

The tiny SSM drafter (27M backbone, $\sim$55\,MB in FP32 excluding the shared embedding) fits comfortably in CPU memory, freeing all GPU memory for the verifier.
We implement an optimized CPU inference stack in three progressively faster tiers.

\paragraph{Tier 1: AVX-512 SSM kernel.}
The core selective scan recurrence (Eq.~\ref{eq:ssm}) processes $N\!=\!128$ state elements per $(h,d)$ position.
With 512-bit SIMD registers holding 16 float32 values, each state update requires exactly 8 vectorized FMA iterations.
We parallelize over the $(B, H, D)\!=\!1{\times}16{\times}64$ independent work items using OpenMP.
This achieves $4\times$ speedup over PyTorch CPU for the SSM kernel alone.
However, profiling reveals that linear projections dominate per-layer time (85\%), not the SSM step (10\%).

\paragraph{Tier 2: Fused BF16 forward pass.}
We fuse the entire 16-layer forward pass and LM head into a single C++ function, eliminating Python dispatch overhead ($\sim$3200 calls per token) and converting linear layers to BF16 with AMX matrix multiply.
The fused BF16 model achieves $2.8\times$ over na"ive PyTorch, reducing 8-token draft latency from 79\,ms to 32\,ms at 8 threads.

\paragraph{Tier 3: INT8 VNNI quantization.}
We quantize weights and activations to INT8 using AVX-512 VNNI instructions, which perform 4 multiply-accumulates per clock cycle on supported hardware.
The INT8 fused model achieves $7.5\times$ over na"ive PyTorch, completing 8-token drafts in 9\,ms at 16 threads (924\,tok/s throughput).

\paragraph{Asynchronous pipeline.}
Since the LLaMA-8B verifier takes $\sim$16\,ms per forward pass on an H100 GPU and our INT8 CPU draft takes $\sim$9\,ms for 8 tokens, the draft fits entirely within the verification window.
In the pipelined regime, the GPU verifies the current batch while the CPU concurrently drafts the next batch.
With mean acceptance of 2.33 tokens per round, the projected throughput is 202\,tok/s, a $2.7\times$ improvement over the 75.8\,tok/s autoregressive baseline.
